\documentclass[10pt,a4paper]{amsart}
\usepackage[margin=19mm]{geometry}
\usepackage{amsmath,amssymb,mathtools}
\usepackage[scr=rsfs,cal=euler]{mathalfa}
\usepackage{xfrac}
\usepackage[unicode,hidelinks,bookmarks=false]{hyperref}
\newcommand{\ie}{i.e.\ }
\newcommand{\cf}{cf.\ }
\newcommand{\resp}{respectively\ }
\newcommand{\Subsection}{\subsection}
\providecommand{\nobreakdash}{\nobreak\hbox{-}}
\setlength{\emergencystretch}{2em}
\multlinegap0pt
\usepackage{amssymb}
\usepackage{mathtools}
\usepackage{xcolor}
\newtheorem{lemma}{Lemma}
\title{A Short Combinatorial Proof of the Pons--Batle Identity for Counting Tree-Child Networks}
\author{Hao Yu, Louxin Zhang}
\date{Source: arXiv:2609.04979v1 (CC BY 4.0)}
\begin{document}
\maketitle

\noindent\emph{Source and licence.} A Short Combinatorial Proof of the Pons--Batle Identity for Counting Tree-Child Networks. Version arXiv:2609.04979v1 (CC BY 4.0), distributed by arXiv under the Creative Commons Attribution 4.0 International licence, http://creativecommons.org/licenses/by/4.0/, as recorded in the arXiv licence metadata for this exact version and shown on the abstract page https://arxiv.org/abs/2609.04979v1. Full licence notice: \texttt{licenses/arxiv-2609.04979-CC-BY-4.0.txt}.\par\smallskip

\noindent\textit{Editorial scope.} Counted unit: the article's lemma lem:twomark (source0.tex lines 678--684). The article's complete original proof of it is delivered with it (source0.tex lines 686--792, one \texttt{proof} environment of 107 lines). No mathematics is written, repaired or re-derived: the statement and the proof are the article's own lines, in the article's own notation, and no retained line is substituted or re-flowed.  No further source-local context is needed: every cross-reference of every retained line resolves inside the two delivered spans.  Nothing was omitted from any retained span, so the package carries no omission record.  No retained line contains a citation, so the wrapper carries no bibliography and no bibliography entry.  No retained line contains a figure environment, an image-inclusion command or drawing code, so no image is delivered and the package declares no asset dependency.  Editorial formatting only, in addition: an AMS \texttt{amsart} preamble that restates the article's own declarations and macros used by the retained lines, namely the environments \texttt{lemma} on separate counters, together with the standard packages those lines need.  The retained environments print in this document as Lemma 1 in order of appearance, whereas the article prints the same items with the numbering of its own full document.  The complete native source of the article as distributed in the arXiv e-print for this exact version is bundled unchanged as \texttt{sources/209401/source0.tex}.
\begin{lemma} \label{lem:twomark}
For $n\ge2$ and $k\ge1$,
\begin{eqnarray}
 2\,D(n,k)=\sum_{\substack{m+p=n\\ i+j+c=k}}
  \Bigl[K_1\,D(m,i)\,a_{p, j}+K_2\,a_{m, i}\,D(p,j)+K_3\,U(m,i)\,a_{p, j}\Bigr]. \label{eq:twomark}
\end{eqnarray}
\end{lemma}

\begin{proof} 
Applying the root-split identity at $(n,k-1)$ gives
\begin{eqnarray*}
2a_{n,k-1}
=
\sum_{\substack{m+p=n\\ \alpha+\beta+r=k-1}}
\binom{n}{m}2^r
\binom{p-\beta}{r}
(2m+\alpha-1)^{\overline r}
a_{m,\alpha}a_{p,\beta}.
%\tag{12}
%\label{eqn99}
\end{eqnarray*}
Hence,
\begin{eqnarray}
2D(n,k)
&={}&
\sum_{\substack{m+p=n\\ \alpha+\beta+r=k-1}}
\binom{n}{m}2^r
\binom{p-\beta}{r}
(2m+\alpha-1)^{\overline r}  \nonumber\\
& &\qquad{}\times
(n-k)(n-k+1)a_{m,\alpha}a_{p,\beta}.
%\tag{13}
\label{eqn133}
\end{eqnarray}

For a summand in (\ref{eqn133}), put $
x=m-\alpha,\; y=p-\beta-r$.
Since $m+p=n$ and $\alpha+\beta+r=k-1$, we have 
$x+y=n-k+1$.
Therefore,
\[
(n-k)(n-k+1)
=(x+y-1)(x+y)
=x(x-1)+y(y-1)+2xy.
\]
Substituting this into (\ref{eqn133}), we obtain

\begin{eqnarray}
2D(n,k)
&={}&
\sum_{\substack{m+p=n\\ \alpha+\beta+r=k-1}}
\binom{n}{m}2^r
\binom{p-\beta}{r}
(2m+\alpha-1)^{\overline r}  \nonumber \\
& &\qquad{}\times
\bigl[x(x-1)+y(y-1)+2xy\bigr]
a_{m,\alpha}a_{p,\beta}.
\label{eqn144}
\end{eqnarray}


We reindex the three terms in (\ref{eqn144}) separately.

For the term $x(x-1)$, set
\[
i=\alpha+1,\qquad j=\beta,\qquad c=r.
\]
Then $i+j+c=k$ and
$D(m,i)=x(x-1)a_{m,\alpha}$.
Moreover, $2m+\alpha-1=2m+i-2=s-1$.
Thus this part of (\ref{eqn144}) becomes
\[
\sum_{\substack{m+p=n\\i+j+c=k}}
K_1D(m,i)a_{p,j}.
\]

For the term $y(y-1)$, set
\[
i=\alpha,\qquad j=\beta+1,\qquad c=r.
\]
Then $i+j+c=k$. Since $t=p-j=p-\beta-1$, we have
$y=p-\beta-r=t+1-c$.
Using
\[
(t+1-c)(t-c)\binom{t+1}{c}
=
t(t+1)\binom{t-1}{c}
\]
and $D(p,j)=t(t+1)a_{p,\beta}$,
this part becomes
\[
\sum_{\substack{m+p=n\\i+j+c=k}}
K_2a_{m,i}D(p,j).
\]

Finally, for the term $2xy$, set
\[
i=\alpha,\qquad j=\beta,\qquad c=r+1.
\]
Then $i+j+c=k$, $x=m-i$, and $y=t-c+1$.
Using
\[
(t-c+1)\binom{t}{c-1}
=
c\binom{t}{c}
\]
together with $2\cdot2^{c-1}=2^c$, this part becomes
\[
\sum_{\substack{m+p=n\\i+j+c=k}}
K_3U(m,i)a_{p,j},
\]
where $U(m,i)=(m-i)a_{m,i}$.

 Consequently, \eqref{eq:twomark} holds.
\end{proof}

\end{document}
